\ulohahead{specializace UI-SU}{k-nejbližších sousedu (kNN)}
\begin{zadani}
\begin{enumerate}[label=(\alph*)]
  \item \bod{1} Popište algoritmus k-nejbližších sousedu pro klasifikaci i regresi. Proč se mu říká ,,líné'' učení (lazy)?
  \item \bod{2} Jak se volí $k$ (křížová validace) a jak malé vs velké $k$ ovlivní výsledek (vychýlení vs rozptyl)? Proč je nutné škálovat příznaky?
  \item \bod{3} Co je prokletí dimenzionality a jak postihuje kNN? Jaká je výpočetní složitost predikce?
\end{enumerate}
\end{zadani}

\begin{reseni}
\textbf{(a)} kNN si \emph{pamatuje} trénovací data; pro nový bod najde $k$ nejbližších (dle metriky, typicky euklidovské) a \emph{klasifikace} = většinové hlasování jejich tříd, \emph{regrese} = průměr jejich hodnot. Je \emph{líné}, protože netrénuje žádný model - veškerá práce je až při predikci.

\textbf{(b)} $k$ se volí \emph{křížovou validací}. \emph{Malé} $k$ (např.\ 1) dává nízké vychýlení, ale vysoký rozptyl (citlivé na šum, přeučení); \emph{velké} $k$ hranici vyhladí (vyšší vychýlení, nižší rozptyl). \emph{Škálování} je nutné, protože vzdálenost jinak ovládne příznak s největším rozsahem.

\textbf{(c)} \emph{Prokletí dimenzionality}: ve vysoké dimenzi se vzdálenosti mezi body vyrovnávají (poměr nejbližší/nejvzdálenější $\to1$), takže pojem ,,soused'' ztrácí smysl a kNN selhává; navíc je třeba exponenciálně víc dat. \emph{Složitost} naivní predikce je $O(nd)$ na jeden dotaz (projít $n$ bodu v dimenzi $d$); urychlení přes k-d stromy či přibližné hledání.
\end{reseni}

\begin{jadro}
kNN = ulož data, pro nový bod najdi $k$ nejbližších; klasifikace = hlasování, regrese = průměr. Líné učení. $k$ přes křížovou validaci (malé $k$ = rozptyl, velké = vychýlení). Škáluj příznaky; trpí prokletím dimenzionality.
\end{jadro}

\kalibrace{kNN (~12\,\%) je vzácnější, ale jednoduchý a vděčný. Definice + role $k$ a škálování = spolehlivá 2-bodová záloha; navíc nese pojem prokletí dimenzionality.}
\past{kNN netrénuje (lazy) - ,,trénink'' je jen uložení dat. Bez škálování dominuje příznak s velkou jednotkou; pozor i na liché $k$ u dvou tříd (remízy).}
